\documentclass[10pt,a4paper]{amsart}
\usepackage[margin=19mm]{geometry}
\usepackage{amsmath,amssymb,mathtools}
\usepackage[scr=rsfs,cal=euler]{mathalfa}
\usepackage{xfrac}
\usepackage[unicode,hidelinks,bookmarks=false]{hyperref}
\newcommand{\ie}{i.e.\ }
\newcommand{\cf}{cf.\ }
\newcommand{\resp}{respectively\ }
\newcommand{\Subsection}{\subsection}
\providecommand{\nobreakdash}{\nobreak\hbox{-}}
\setlength{\emergencystretch}{2em}
\multlinegap0pt
\usepackage{bm}
\usepackage{bbm}
\usepackage{dsfont}
\usepackage{xspace}
\usepackage{cancel}
\usepackage{mathtools}
\usepackage[shortlabels]{enumitem}
\usepackage{amsthm}
\makeatletter
\@ifundefined{lem}{\newtheorem{lem}{Lemma}}{}
\makeatother
\title[Relaxed Newton's Method as a Family of Root-finding Methods:]{Relaxed Newton's Method as a Family of Root-finding Methods: Dynamics and Convergence}
\author{Soumen Pal}
\date{Source: arXiv:2603.11591v1 (CC BY 4.0)}
\begin{document}
\maketitle

\noindent\emph{Source and licence.} Relaxed Newton's Method as a Family of Root-finding Methods: Dynamics and Convergence. Version arXiv:2603.11591v1 (CC BY 4.0), distributed under the Creative Commons Attribution 4.0 International licence, as recorded in the arXiv licence metadata for this exact version and shown on the abstract page https://arxiv.org/abs/2603.11591v1. Full rights notice: licenses/arxiv-2603.11591-CC-BY-4.0.txt.\par\smallskip

\noindent\textit{Editorial scope.} Counted unit: the Lem whose source label is \texttt{None} in the article (source0.tex lines 239--249, source label \texttt{None}), delivered together with the article's own complete original proof of it (source0.tex lines 250--282, one \texttt{proof} environment of 33 lines). No mathematics is written, repaired or re-derived: statement and proof are the article's own text line for line. Required context retained uncounted: FPI\_simple (lines 151--153). Every definition, display and label of those lines is retained unchanged. Omitted from this package and not redistributed: 1--150, 154--238, 283--619. Nothing inside a retained excerpt is omitted or altered: every retained body is delivered exactly as the article's lines are written, so no omission receipt of any kind accompanies this package. Citations: the retained excerpts carry no citation command at all, so no bibliography entry of the article is transcribed and no reference is redistributed. Reference adaptation: none was needed and none was made; every reference inside the delivered unit resolves inside it, so no unresolved reference or citation is produced. Printed slips kept literal: no printed word, symbol, hypothesis or number of the counted statement or of its proof was altered. Layout and class: the article's own class is \texttt{article} with packages including bbm, graphicx, wrapfig, euscript, subcaption, hyperref, epstopdf, enumerate, longtable, float, amsmath, amsfonts, amssymb, amsthm; the wrapper is the lane's pinned \texttt{amsart} editorial wrapper at 10pt on A4, which restates the theorem-like environments the retained text needs and adds the article's own macro definitions that the retained text needs (none), with extra wrapper packages (bm, bbm, dsfont, xspace, cancel, mathtools, enumitem). The delivered numbering is the unit's own. Assets: no retained span contains a figure environment, a graphics-inclusion command, a PDF-page inclusion, a table or a TikZ picture; no image file is copied into this package, the package declares no asset dependency and no image file is redistributed with it. Licence: version arXiv:2603.11591v1 of this article is distributed under the Creative Commons Attribution 4.0 International licence, as recorded in the arXiv licence metadata for this exact version and shown on the abstract page https://arxiv.org/abs/2603.11591v1. A full rights notice is delivered at licenses/arxiv-2603.11591-CC-BY-4.0.txt. Characters: every retained excerpt is pure ASCII, so the delivered engine, XeLaTeX with its default Latin Modern text font, reports no missing character.
\begin{equation}\label{FPI_simple}
\iota (R,z_0)=\frac{1}{1-\lambda}.
\end{equation}

\begin{lem}[Characterization of quadratic relaxed Newton maps]
Let $R$ be a quadratic rational function having two attracting and one repelling fixed point. Then the following are true.
\begin{enumerate}[(i)]
\item If the multipliers of two attracting fixed points are the same, say $\lambda$, then $R$ is conjugate to $N_{h, p}$, where
$
h=1-\lambda \text { and } p(z)=z^{2}-1
$.
\item If one of the fixed points of $R$ is superattracting and the multiplier of another attracting fixed point is $\frac{n}{m}$, where $m\in \mathbb{N}$ and $n \in \mathbb{N}\setminus \{0\}$, then $R$ is conjugate to $N_{h, p}$, where $h=m-n$ and $p(z)=(z-1)^{m-n}(z+1)^{m}$.
\item If the ratio of the residues of the fixed point indices corresponding to the attracting fixed points of $R$ is $\frac{k}{m}$, where $k,m\in \mathbb{N}$, then there exists $h\in \mathbb{C}\setminus \{0\}$ such that $R$ is conjugate to $N_{h, p}$, where $p(z)=(z-1)^{k}(z+1)^{m}$.
 \end{enumerate}
\end{lem}

\begin{proof}
Let $R$ be a quadratic rational function having three distinct fixed points, among them, one being repelling and the other two being attracting. Using conjugacy, we can consider $ 1 $ and $ -1 $ to be attracting, whereas $\infty$ is the repelling fixed point of $R$. Thus, $R$ has the following form:
\begin{equation}\label{form_R}
R(z)=z-\frac{z^{2}-1}{A z+B},
\end{equation}
where $A, B \in \mathbb{C}$. Then
\begin{equation}\label{mult_R}
R^{\prime}(1)=1-\frac{2}{A+B} \text { and } R^{\prime}(-1)=1-\frac{2}{A-B}.
\end{equation}
\begin{enumerate}
	\item If $R^{\prime}(1)=R^{\prime}(-1)=\lambda$ then $B=0$ and $\lambda=1-\frac{2}{A}$.
Therefore from (\ref{form_R}), we have
$$R(z) =z-\frac{z^{2}-1}{A z}=z-\frac{2}{A} \frac{z^{2}-1}{2 z} \\
=N_{h, p}(z),$$
where $p(z)=\left(z^{2}-1\right)$ and $h=\frac{2}{A}=1-\lambda$.
\item  Let $R^{\prime}(1)=0$
and $R^{\prime}(-1)=\frac{n}{m}$. Then we get $A+B=2$ and $A-B=\frac{2 m}{m-n}$.
Therefore, $A=\frac{2 m-n}{m-n}$ and $B=-\frac{n}{m-n}$.
Hence from (\ref{form_R}),
$$
\begin{aligned}
R(z) & =z-\frac{(m-n)\left(z^{2}-1\right)}{(2 m-n) z-n} \\
& =z-(m-n) \frac{z^{2}-1}{((m-n)+m) z+((m-n)-m)} =N_{h, p}(z),
\end{aligned}
$$
where $h=m-n$ and $p(z)=(z-1)^{m-n}(z+1)^{m}$.
\item  From (\ref{FPI_simple}), we have $\iota(1)=\frac{1}{1-R'(1)}$ and $\iota(-1)=\frac{1}{1-R'(-1)}$. Thus $\frac{\iota(1)}{\iota(-1)}=\frac{1-R^{\prime}(-1)}{1-R^{\prime}(1)}=\frac{k}{m}$. Following from part two of this proof, without loss of generality we can assume that $1<k<m$. Now it follows from (\ref{mult_R}) that
$\frac{A+B}{A-B}=\frac{k}{m}.$
Therefore, there exists $c \in \mathbb{C} \setminus \{0\}$ such that $A+B=2 k c$ and $A-B=2 m c$. Thus we get $A=(m+k) c$ and $B=(m-k) c$. This gives that $$R(z)=z-\frac{1}{c} \frac{z^{2}-1}{(k+m) z+(k-m)}=N_{h, p},$$
 where $h=\frac{1}{c}$ and $p(z)=(z-1)^{k}(z+1)^{m}$.
\end{enumerate}
This concludes the proof.
\end{proof}

\end{document}
